% Auto-generated from 41-Transition-Decision.md
% POV: Aisen | Landscape: Dungeon | Location: The Covenant Halls

Aisen spent the final month of military operational phase while maintaining simultaneous coordination through civilian resistance network channels. The transition decision was approaching: would he remain positioned in military command structure to coordinate during institutional replacement phase, or would he exit military employment to work directly with provincial resistance networks?

The advantages of remaining in military position were significant: access to classified institutional planning, authority to coordinate military personnel attachments to civilian governance structures, visibility into Covenant Council decision-making about administrative implementation. The advantages of transitioning to direct resistance work were equally significant: operational independence from military authority, ability to coordinate provincial resistance networks directly, positioning to influence institutional replacement from civil society perspective rather than from internal military position.

Kael and Aisen met one final time in Southeastern Territory to discuss the decision.

"If you stay in military structure, you can help us understand institutional replacement implementation," Kael said. "That intelligence would be valuable during the transition phase."

"And if I leave?" Aisen asked.

"Then you can work directly on institutional subversion and resistance network coordination," Kael said. "You can help us prepare alternative governance frameworks. You can position provincial networks to challenge Covenant administrative control as soon as transition phase concludes."

"Which is more strategically valuable?" Aisen asked.

"Uncertain," Kael admitted. "Both approaches have merit. My assessment is that institutional subversion probably matters more long-term, but institutional intelligence definitely matters more short-term during transition implementation."

"What does Amara recommend?" Aisen asked.

"She recommends you remain in military structure," Kael said. "She thinks long-term military intelligence asset is more valuable than short-term resistance network coordination."

Aisen made his decision.

"I'm going to transfer to provincial resistance network operations," he said. "After formal military occupation maintenance phase begins and military forces transition, I'll exit military employment and transition to full-time resistance work."

The military campaign formally concluded with Covenant Council authorization for institutional replacement implementation. Military forces transitioned to garrison positions in major urban centers. Covenant administrative personnel began arriving in provincial territories to establish governance structures.

Aisen submitted his military discharge request citing "family medical circumstances requiring return to civilian status." The request was administratively processed through normal military channels and approved within two weeks. He was officially separated from military service with rank of Captain and classified as eligible for recall if military emergency circumstances required.

He left military installations and returned to civilian status in Southeastern Territory with official permission to remain in the region under Covenant civilian governance authority.
